\scanpage{6}

% ll. 1–3
\mspage{\textcircled{4}}
Il faut montrer que le $2^{\mathrm{ième}}$ mb ne dépend que de la partition $\in\Pi_{22}(P_u)$, pas du choix de $(x,y)$. Pour ceci, on note

% ll. 4–5
a) Elle est indépendante par permutation de $x$ et $y$ \deleted{\unsure{c'est}} (car $\delta_v(x,y)=\delta_v(y,x)$)

% ll. 6–8
b) Elle ne change pas si on change $x$ et $\sigma(x)$ (ou aussi $y$ en $\sigma(y)$, ou $x$ en $\sigma(x)$ et $y$ en $\sigma(y)$) (car $\delta_v(\sigma x,y)=\delta_v(x,y)+1$)

% ll. 9–13
c) Si $x,y$ sont fixés, \deleted{\unsure{prenons|posons}}\inserted{posant (lemme 4)} \deleted{$s\ne t$} $\{s,t\}=\deleted{E(x)\cap E(y)\cap E(v)^*}E(x)\cap E(y)\cap E(v)^*$, donc \inserted{(id)} $\{x,y\}=E(s)\cap E(t)\cap E(u)^*$, \deleted{soient}\inserted{soit} $\{\unsure{\xi|z|x'},\eta\}=E(s)\cap E(\sigma(t))\cap E(u)^*$, alors $\pi_u(\{\xi,\eta\})=\complement_{P_u}\pi_u(\{x,y\})$, \deleted{\unsure{donc $x',y$|donc $x',y'$}} et $\{\xi,\eta\}$ donne donc lieu à la $=$ partition $\in\Pi_{2,2}(P_u)$ que $\{x,y\}$,

% ll. 14–19
mais à $\{\xi,\eta\}$ la recette \deleted{\unsure{préc.|prise}} de l'énoncé du corollaire associe la partition en $\pi_v(\underbrace{E(\xi)\cap E(\eta)\cap E(v)^*}_{\{s,\sigma(t)\}})$ et son complé\unsure{mentaire|mtaire}, qui \deleted{\unsure{sont}} égaux à $\underbrace{\pi_v\{s,t\}}_{\operatorname{supp}(\delta_v(x,y)+1)}$ et son complé\unsure{mentaire|mtaire} \unsure{---|\dots}

% l. 20
Énoncé plus fort :

% ll. 21–22
\deleted{\unsure{Cor}}\uline{Corollaire 2} Soit $\mathcal P'_2(E_u^*)$ l'ens des \deleted{couples}\inserted{parties à} 2 éléments de $E_u^*$ formés de deux él. non \uline{liés}.

% ll. 23–25
Soit $R_{u,v}$ la relation \deleted{entre} dans $\mathcal P'_2(E_u^*)\times\mathcal P'_2(E_v^*)$ : $R_{u,v}(\deleted{\unsure{\beta}}\alpha,\beta)\overset{\text{déf}}{\Longleftrightarrow}$ tt el. de $\alpha$ est lié à tt el. de $\beta$. Alors $R_{u,v}$ \inserted{définit} une bijection

% lower left, oblique
% slanted
\margin{\unsure{énoncé|énoncés|\deleted{Énoncé}} plus simple avec $R'_{uv}(\alpha,\beta)$ : tt él. de $\alpha$ \uline{non lié} à tt él. de $\beta$, $[\Longleftrightarrow R_{uv}(\sigma\alpha,\beta)\Longleftrightarrow R_{uv}(\alpha,\sigma\beta)]$ car les \unsure{applications|appl:ations} $\varphi'_{vu}:\mathcal P_2(E_u^*)\to\mathcal P_2(E_v^*)$ forment un syst. \uline{transitif} d'iso (comme on \unsure{le verra|verra} au lemme 6 ci-dessous)}

% ll. 26–27
de $\mathcal P'_2(E_u^*)$ sur $\mathcal P'_2(E_v^*)$, soit \deleted{$\varphi$}

% ll. 27–29
\[
\begin{aligned}
\widetilde\varphi_{vu}:\mathcal P'_2(E_u^*)&\longrightarrow\mathcal P'_2(E_v^*)\\
\widetilde\varphi_{vu}(\sigma\widetilde\alpha)&=\sigma\widetilde\varphi_{vu}(\widetilde\alpha)\\
\pi_v\widetilde\varphi_{vu}(\{x,\sigma y\})&=\complement_{P_v}\pi_v\varphi_{vu}(\{x,y\})
\end{aligned}
\]
% archive sheet number: 5
